% !TeX root = ../main.tex
\chapter{Background and Related Work}\label{ch:background}

\begin{guidance}
Show that you know the field. First describe the problem and the concepts and techniques
needed to understand the thesis (theoretical framework); then review critically the most
relevant work published on the topic (state of the art), grouping it by approach rather
than listing papers one by one. End by identifying the gap that your work addresses.
References are essential here: search the scientific literature (Scopus, Web of Science,
Google Scholar, PubMed, IEEE Xplore, ACM Digital Library\dots) and cite every source.
\end{guidance}

\section{Theoretical framework}\label{sec:framework}

[Introduce the concepts, definitions and techniques on which the thesis is built. For
example, an introduction to \gls{ml} for multi-label problems can be found
in~\cite{herrera2016multilabel}.]

\section{State of the art}\label{sec:state-of-the-art}

[Review and compare the most relevant previous works. A summary table that compares the
approaches along a few dimensions is often very useful.]

\section{Summary and research gap}\label{sec:gap}

[Summarise what is known and, most importantly, what is not known yet or not solved,
which motivates the objectives in \cref{ch:objectives}.]
